\documentclass[11pt,a4paper]{article}
\usepackage[T1]{fontenc}
\usepackage{lmodern,amsmath,amssymb,microtype,booktabs}
\usepackage[margin=25mm]{geometry}
\usepackage[colorlinks=true,urlcolor=blue,linkcolor=blue]{hyperref}
\newcommand{\dd}{\mathop{}\!\mathrm d}
\newcommand{\cJ}{\mathcal J}
\title{Crossed recoupling in the local dS$_3$ model:\\physical normalization and a Euclidean check}
\author{Alan Day \and Chengheng\thanks{AI research assistant (OpenAI).}}
\date{Working research note --- 31 October 2026}
\begin{document}\maketitle
\begin{abstract}
We match the scalar channels reached by two ordered local creations to
standard conformal partial waves. Three-point group normalization and an
independent partial-wave norm fix the conversion to the known conformal
$6j$ symbol, including its phase. A direct Euclidean area integral checks
the complex matrix element independently of its Lorentzian inversion
formula. We give a conservative analytic bound on the three puncture
regions and formulate the crossed overlap for physical creation states.
The calculation concerns a declared $G_N=0$ model. The scalar block is a
compression of the full Racah transform; it is not claimed to be unitary.
The full physical wave-packet integral remains to be evaluated.
\end{abstract}

\section{Model, channels, and normalization}
Set the de Sitter radius to one, $G_N=0$, and use the proper group
$SO^+(3,1)$. Two scalar Bunch--Davies observers have labels at least
$s_*>0$. The distinct free matter species have positive labels $p,q$;
local couplings are one. Matter modes have Klein--Gordon normalization,
and observer modes have the previously fixed spacetime delta normalization.
Write
\[
 h_{\rm m}(x)=\sqrt{1+x^2},\qquad \rho(x)=2x^2/\pi,\qquad
 \mu_{\rm obs}(x)=x\tanh(\pi x)/\pi .
\]
The compact subgroup has Haar probability, with radial push-forward
$\sinh^2 r\,\dd r$. For $S=x+y+z$, retain
\[
 D_{xyz}=\cosh(\pi S/2)\prod_{n=x,y,z}\cosh[\pi(S-2n)/2],\qquad
 K_{xyz}=\frac{\pi}{8D_{xyz}}\prod_{n=x,y,z}\frac{\sinh(\pi n)}n .
\]
These conventions and the physical four-body norms were derived in our
preceding two-matter note. They are inputs here.

Fix the external order $e=(a,q,c,p)$. The $s$ pairing is
$(a,q)|(c,p)$, while the LPRS $t$ pairing is $(c,q)|(a,p)$
\cite{lprs}. The latter is the reverse of the pair order naturally used
by the opposite creation sequence. This reversal must be accompanied by
an internal shadow before normalization.
The full CG decomposition includes all integer principal-series spins
\cite{ds,kks}. Each local scalar observer creation reaches spin zero in
its own pairing. This does not imply that the scalar sector is complete
after recoupling. We calculate the $00$ matrix element of the full transform.

\section{From sphere periods to physical channel states}
Let $\Delta_x=1+ix$. Denote the all-ket boundary trilinear's spherical
$000$ component by $B(x,y,z)$. For $\nu_i=(x,y,z)$, put
\[
 \alpha_i=-\frac12-\frac i2(S-2\nu_i),\quad A=\sum_i\alpha_i,\qquad
 B=(4\pi)^{3/2}4^A
 \frac{\Gamma(A+2)\prod_i\Gamma(\alpha_i+1)}
 {\prod_i\Gamma(A-\alpha_i+2)} .
\]
Its modulus is
\begin{equation}
 |B(x,y,z)|^2=\frac{\pi^4}{D_{xyz}}
                 \prod_{n=x,y,z}\frac{\sinh(\pi n)}n .
 \label{Bnorm}
\end{equation}
Negative labels in $B$ denote shadow representatives, not negative masses.
Use standard shadow-integral partial waves
\[
 \Psi_b^s=\int\dd^2x\,
  t_{\Delta_a\Delta_q\Delta_b}\,
  t_{\widetilde\Delta_b\Delta_c\Delta_p},
 \qquad
 \Psi_d^t=\int\dd^2x\,
  t_{\Delta_c\Delta_q\Delta_d}\,
  t_{\widetilde\Delta_d\Delta_a\Delta_p}.
\]
Projection of two external spherical $l=0$ states leaves only internal
$l=0$. The remaining normalized spherical harmonics pair to one.
Consequently the $0000$ components are exactly
\begin{equation}
 N_s(b)=B(a,q,b)B(c,p,-b),\qquad
 N_t(d)=B(c,q,d)B(a,p,-d).
 \label{Ns}
\end{equation}
The plane-to-sphere unitary conformal map introduces no further constant:
the two contracted dimensions sum to two and absorb its area Jacobian.
In the unique scalar CG channel, the $0000=1$ states therefore are
\begin{equation}
 J_s(b)=\Psi_b^s/N_s(b),\qquad J_t(d)=\Psi_d^t/N_t(d).
 \label{Jnorm}
\end{equation}
For conformal weights $h=(1+ib)/2$, $h_i=(1+ie_i)/2$, use
\[
 K_h^{ij}=\pi\frac{\Gamma(2h-1)}{\Gamma(2-2h)}
 \frac{\Gamma(1-h-h_i+h_j)\Gamma(1-h+h_i-h_j)}
 {\Gamma(h-h_i+h_j)\Gamma(h+h_i-h_j)} .
\]
The Gamma identities in $B$ and the two shadow-block coefficients give
\begin{equation}
 \frac{N_s(b)}{N_s(-b)}=\frac{K_{1-h}^{34}}{K_{1-h}^{12}},
 \qquad
 \frac{\Psi_b^s}{N_s(b)}=\frac{\Psi_{-b}^s}{N_s(-b)} .
 \label{reflection}
\end{equation}
Thus the raw waves need not be equal under reflection. Exchanging the
two trilinears also reflects the internal label; after division by
$N$, it gives the same $J$. This connects the LPRS $t$ orientation to
the actual $(a,p)|(c,q)$ creation channel.

\section{The physical crossed kernel}
KKS fixes the scalar three-point quotient pairing to
$[2^d\operatorname{vol}SO(d-1)]^{-1}$, hence $1/4$ at $d=2$
\cite[(13)--(15)]{kks}. In our group convention,
\[
 \|t_{xyz}\|_{\rm phys}^2=|B(x,y,z)|^2/K_{xyz}=8\pi^3 .
\]
The common Haar-volume ratio between our physical quotient and the
LPRS/KKS quotient is therefore $32\pi^3$. It is fixed by the
three-point pairing, not fitted to a four-point value.
The previously derived scalar four-body rigging weights are
\[
 W_s(b)=\rho(b)K_{abq}K_{bcp},\qquad
 W_t(d)=\rho(d)K_{cdq}K_{dap}.
\]
A distribution $\int F(b)J_s(b)\,\dd b$ has physical norm
$\int |F(b)|^2/W_s(b)\,\dd b$. Equation \eqref{Bnorm} gives
\begin{equation}
 \frac{|N_s(b)|^2}{W_s(b)}=\frac{32\pi^7}{b^2}.
 \label{rawn}
\end{equation}
Independently, the scalar shadow square is $\mathcal N_b=\pi^2/b^2$.
KKS (74) gives $\mu_{\rm KKS}=b^2/(8\pi^3)$, distinct from
$\mu_{\rm obs}$. Their (39) yields, with positive-label measure $\dd b$,
\[
 n_b=(1/4)^2\,2\pi/\mu_{\rm KKS}=\pi^4/b^2.
\]
This agrees with \eqref{rawn} divided by $32\pi^3$.
One reflection representative has already been selected; no extra
factor of two is present.
Let $\cJ_2(b,d;e)=\langle\Psi_b^s,\Psi_d^t\rangle_{\rm LPRS}$
be the known scalar conformal $6j$ symbol \cite[(3.4)--(3.5)]{lprs}.
For delta-normalized states $e_b^s=\sqrt{W_s(b)}J_s(b)$,
\begin{align}
 \langle J_s(b),J_t(d)\rangle_{\rm phys}
 &=\frac{32\pi^3\cJ_2(b,d;e)}{\overline{N_s(b)}N_t(d)},\\
 U_{00}(b,d;e)&=\frac{bd}{\pi^4}
   \frac{N_s(b)/|N_s(b)|}{N_t(d)/|N_t(d)|}\,\cJ_2(b,d;e).
 \label{U}
\end{align}
The phases cannot be replaced by modulus squares.
This dictionary follows the actual local CG chains. It does not undo
the previously found normalization-invariant mismatch of the separate
$H/Z$ candidate kernels.

\section{An independent Euclidean overlap}
Gauge-fix $(z_1,z_2,z_3,z_4)=(0,z,1,\infty)$ and define
\[
 k_h(z)=z^h\,{}_2F_1(h-h_{12},h+h_{34};2h;z).
\]
For scalar principal series, $\bar h=h$, not $h^*$. The full wave is
\begin{equation}
 F_s(z)=|z|^{-2(h_1+h_2)}
 \left[K_{1-h}^{34} k_h(z)k_h(\bar z)
       +K_h^{12}k_{1-h}(z)k_{1-h}(\bar z)\right].
 \label{Fs}
\end{equation}
For $F_t$, replace the external order by $(3,2,1,4)$, the internal
label by $d$, and $z$ by $1-z$. The direct check is
\begin{equation}
 \cJ_2=\frac14\int_{\mathbb C}\dd^2z\,
                      \overline{F_s(z)}F_t(z).
 \label{euclidean}
\end{equation}
This never calls the Lorentzian inversion's $ {}_4F_3$ or Euler-product
evaluator. It shares only the standard three-point shadow coefficients.
Use the base domain $|z|<1$, $|1-z|>1$ and its six images
\[
 z,\quad 1-z,\quad 1/z,\quad 1/(1-z),\quad z/(z-1),\quad (z-1)/z.
\]
They cover the plane once away from boundaries: their domains correspond
to the six orderings of $|z|,|1-z|,1$. Their area Jacobians are
$1,1,|z|^{-4},|1-z|^{-4},|1-z|^{-4},|z|^{-4}$.
On the upper half put $z=r e^{i\theta}$ with
$\arccos(r/2)<\theta<\pi$. Scalar waves are unchanged by interchanging
$z,\bar z$, so the lower half supplies a factor two. The combined
quotient factor is $1/2$. We use $r=t^P$, $P=4$ or $8$, and
Gauss--Legendre nodes in $t$ and the mapped angle, retaining every
complex chart contribution.

\section{A conservative analytic puncture bound}
The shadow integral has real part one in every dimension. Hence
\[
 |F_s(z)|\le\frac{I(z)}{|z|},\quad
 |F_t(z)|\le\frac{I(z)}{|1-z|},\qquad
 I(z)=\int_{\mathbb C}\frac{\dd^2x}{|x||x-z||x-1|}.
\]
For $r=|z|\le1/4$, split the $x$ plane at radii $2r,1/2,2$.
The contributions obey
$40\pi$, $8\pi\log(1/(4r))$, $48\pi$, $16\pi/7$, respectively, so
\begin{equation}
 I(z)\le 8\pi[12+\log(1/r)].
 \label{Ibound}
\end{equation}
For the first region set $x=ry$ and use
$|y|+|y-z/r|\ge1$ to bound the reciprocal product by
$2(|y|^{-1}+|y-z/r|^{-1})$. Integrating over disks of radii two and
three gives $40\pi$ after bounding $|x-1|^{-1}\le2$.
In the middle annulus the integrand is at most $4/|x|^2$.
In the third it is at most $8/|x-1|$ on a disk of radius three;
in the last it is at most $16/(7|x|^3)$.

Reflection $z\mapsto1-z$ controls the second puncture. Scaling $x=zy$
gives $I(z)=|z|^{-1}I(1/z)$ for the third. If all six charts omit
$r<\epsilon$, their images lie in disks of radius
$\delta=\epsilon/(1-\epsilon)$ around zero and one, and outside
$|z|=1/\delta$. For $\delta\le1/4$, put $V=12+\log(1/\delta)$
and $C=128\pi^3/3$. Integration gives the raw $6j$ tail bound
\begin{equation}
 T(\epsilon)=C\left[
  2\delta(V^2+2V+2)+
  \frac{\delta^2}{2}\left(V^2+V+\frac12\right)\right].
 \label{tail}
\end{equation}
This proves absolute convergence of the crossed overlap and removal of
the puncture cutoff. It is deliberately conservative and is not a
Gauss quadrature error bound. The same-channel delta norm is a different,
distributional question.

\section{Physical ordered states and the next integral}
Let the initial vacuum coefficient be $f(c)$ with
$\|f\|_0^2=\int\rho(c)|f(c)|^2\dd c$.
Write $E_{xyk}$ for the previously fixed real local $000$ emission
amplitude. At real times, the ordered channel amplitudes are
\begin{align}
 v_s(a,b,c)&=
 \frac{f(c)E_{abq}E_{bcp}}{\sqrt{W_s(b)}}\,
 e^{-it_p[h_{\rm m}(b)-h_{\rm m}(c)]
   -it_q[h_{\rm m}(a)-h_{\rm m}(b)]},\\
 v_t(a,d,c)&=
 \frac{f(c)E_{adp}E_{dcq}}{\sqrt{W_t(d)}}\,
 e^{-iu_q[h_{\rm m}(d)-h_{\rm m}(c)]
   -iu_p[h_{\rm m}(a)-h_{\rm m}(d)]}.
\end{align}
All observer labels are at least $s_*$. Their crossed overlap is
\begin{equation}
 X=\int\dd a\,\dd c\,\dd b\,\dd d\,
      \overline{v_s(a,b,c)}\,U_{00}(b,d;e)\,v_t(a,d,c),
 \label{X}
\end{equation}
understood as the Hilbert-space limit of compact spectral restrictions,
without assuming absolute convergence of the fourfold integral.
The previously proved second-mass-moment condition
$\int\rho(c)|f(c)|^2c^2\dd c<\infty$ is sufficient for the unregulated
two-step states. Shifting both creations by the proper-mass half period
multiplies either amplitude by
$\exp\{\pi[h_{\rm m}(a)-h_{\rm m}(c)]/2\}$; those blocks are bounded
on every initial $L^2$ state.

The full Racah transform is unitary, so its scalar compression is a
contraction. Thus $|X|\le\|v_s\|\|v_t\|$. If compact restrictions give
$v_{s,L},v_{t,L}$, then
\begin{equation}
 |X-X_L|\le
 \|v_s-v_{s,L}\|\|v_t\|+
 \|v_{s,L}\|\|v_t-v_{t,L}\|.
 \label{cutoff}
\end{equation}
These tails are positive norms, computed from the known density
$\rho(c)|f(c)|^2R_q(a|b)R_p(b|c)$ and its exchanged counterpart
(or the half-period rates $\widetilde R$).
This provides a route to a controlled physical packet calculation
without asserting scalar unitarity.

\section{Evidence and limits}
Three label sets pass 33 normalization, phase, channel-exchange, and
shadow-reflection checks at 45 decimal digits. For
\[
 e=(0.4,0.7,0.93,0.6),\quad b=0.65,\quad d=0.57,
\]
the raw and physical matrix elements are
\[
 \cJ_2\simeq190.9971396918+3.7321669491i,\qquad
 U_{00}\simeq0.7266051675149.
\]
The complex Euclidean relative differences at orders $8,16,24,40$
with $P=4$ are approximately
$1.13\,10^{-2}$, $8.93\,10^{-5}$, $5.13\,10^{-6}$, and
$1.01\,10^{-7}$. With $P=8$, orders 24 and 40 give
$2.81\,10^{-7}$ and $8.73\,10^{-11}$.
Stored data retain imaginary parts, chart contributions, absolute
contribution sums, grids, precision, and dependency versions.
For $e=(1.1,0.43,0.82,1.3)$, $b=0.37$, $d=1.17$, the $P=8$
relative differences at orders 24 and 40 are $3.37\,10^{-5}$ and
$4.58\,10^{-8}$. The absolute-sum cancellation ratios are about 2.17
and 3.66 in the two final grids. At the first point, cutting the chart
radius at $10^{-3}$ and $10^{-4}$ gives relative differences
$5.06\,10^{-4}$ and $9.56\,10^{-5}$; these include quadrature error
and are not isolated tail measurements.
The analytic tail bound at $\epsilon=10^{-12}$ is less than
$4.38\,10^{-6}$ in absolute raw-kernel units.
Ten further diagnostics check the chart coverage, the known area
integral $\frac14\int(1+|z|^2)^{-2}\dd^2z=\pi/4$, and the tail
primitive. The area relative error at 32 nodes is below $1.3\,10^{-40}$.
An initial source comment misstated the base domain; its angular
implementation was correct. The executed source is retained, and an
AST comparison confirms that only the comment changed.
These are empirical checks, not interval certificates or uniform
statements over the spectrum.

The new work is the physical normalization dictionary, its coordinate
check, and the explicit domain and cutoff formulation.
The conformal $6j$, CG completeness, and shadow coefficients have the
cited predecessors. The actual fourfold physical packet integral,
identical-species Bose interference, full Fock domains, finite-$G_N$
dynamics, and a complete de Sitter holographic dual remain open.

\begin{thebibliography}{9}
\bibitem{lprs} J.~Liu, E.~Perlmutter, V.~Rosenhaus, D.~Simmons-Duffin,
\emph{$d$-dimensional SYK, AdS Loops, and $6j$ Symbols},
\href{https://arxiv.org/abs/1808.00612v3}{arXiv:1808.00612v3}.
\bibitem{kks} D.~Karateev, P.~Kravchuk, D.~Simmons-Duffin,
\emph{Harmonic Analysis and Mean Field Theory},
\href{https://arxiv.org/abs/1809.05111v1}{arXiv:1809.05111v1}.
\bibitem{ds} S.~E.~Derkachov, V.~P.~Spiridonov,
\emph{The $6j$-symbols for the $SL(2,\mathbb C)$ group},
\href{https://arxiv.org/abs/1711.07073v2}{arXiv:1711.07073v2}.
\end{thebibliography}
\end{document}